\SourceLeaf{306}
\Needspace{5\baselineskip}\begin{center}
[2.]
\end{center}

Man kann, wenn $B, C$ zwei Punkte in $a$ sind, ganz füglich $B - C = ai$ setzen.

\Needspace{5\baselineskip}\begin{center}
[3.]
\end{center}

Es fragt sich nun hauptsächlich, was für Kräfte man, barycentrisch auf die Winkelpunkte, oder resultantisch auf die Seiten des Pentagons appliciren müsse, damit zwischen ihnen und ihren Äquivalenten am innern Pentagon möglichst einfache Verhältnisse Statt finden.

\Needspace{5\baselineskip}\begin{center}
[4.]
\end{center}

Die Collinearität zweier Systeme von Punkten oder geraden Linien ist wohl am treffendsten dadurch zu definiren, dass immer zwischen je 5 Punkten (Linien) des einen Systems und den 5 correspondirenden des andern sich einerlei Lineargleichung aufstellen lässt.

\Needspace{5\baselineskip}\begin{center}
[5.]
\end{center}

Das Characteristische der Collineation scheint am einfachsten so ausgedrückt werden zu können:

Es seien $A, B.\,.\,.$ gerade Linien, $P,Q.\,.\,.$ Punkte des einen Systems, denen im andern die Linien $a,b.\,.\,.$ und die Punkte $p,q.\,.\,.$ entsprechen. Drückt man nun durch $PA$ den (senkrechten) Abstand des Punktes $P$ von der geraden Linie $A$ aus, so wird allgemein
\[
\frac{PA.QB.pb.qa}{pa.qb.PB.QA} = 1.
\]
